\chapter*{Abstract}
\addcontentsline{toc}{chapter}{Abstract}

India presents a paradox. Adult account ownership rose from 35 per cent in 2011 to roughly 89 per
cent by 2025, yet India records the largest share of \emph{inactive} accounts in the world: about
35 per cent of Indian adults who hold an account hold a dormant one, some seven times the average
across other developing economies \citep{findex2021,findex2025}. Accounts have been opened far
faster than they have been used. This report asks whether the expansion of financial infrastructure
is itself part of that story --- whether expansion has been formal rather than functional.

We build verified state-level panels from nine sources: eight Reserve Bank of India and Government
of India series covering branches, ATMs, deposit and credit account counts, deposit and credit
balances, per-capita state income and adult population, plus state-level digital payment data from
PhonePe Pulse. All-India aggregates reconcile to published figures to within 0.001 per cent.
Hypothesis~1 --- that infrastructure raises \Access{} more than \Usage{} --- is tested on three
tracks by two-way fixed-effects regressions of \Access{}, \Usage{} and their difference, with
standard errors clustered by state. The decision rule for accepting or rejecting it was fixed in
writing before any coefficient was estimated.

On the two \textbf{physical} tracks --- 33 states over \panelA{} and 32 over \panelB{} ---
\textbf{the hypothesis is not supported, and the null is uninformative}. A diagnostic run before
estimation, with thresholds fixed in advance, found that only 7.0 per cent of the variation in
branches per lakh adults and 8.7 per cent of that in ATMs occurs \emph{within} states; roughly 93
per cent is cross-sectional, which is exactly what fixed effects discard. A Hausman test then
rejects random effects decisively, ruling out the estimator that would recover that variation. The
fragility this implies is demonstrated rather than asserted: dropping two observations out of 198
reverses the sign of every ATM coefficient.

On the \textbf{digital} track --- 33 states over FY2019--FY2023, adding UPI adoption --- the
identification problem disappears. Within-state variation in UPI users per adult is 67.1 per cent,
and explanatory power rises from a within-$R^2$ of 0.069 to 0.609 for \Access{}. UPI adoption
raises \Access{} ($+0.439$) more than \Usage{} ($+0.183$), and the gap is significant
($p = 0.020$) --- the first time in this study the access--usage gap is distinguishable from zero.
Decomposing the outcome shows the effect runs \emph{entirely through credit accounts, not deposit
accounts}, which is the reverse of what an accounting artefact would produce, since deposit
accounts are the component definitionally entangled with UPI registration. The result survives
dropping any single state.

Two caveats bound this. The pre-committed hypothesis named branches and ATMs and it failed; UPI was
added \emph{after} seeing that null, so the digital result is exploratory by construction --- a
hypothesis generated by this data, not one tested and passed. And while the identification problem
is solved, the causal one is not: fixed effects do not remove state-specific time-varying
confounders.

\vspace{0.8em}
\noindent\textbf{Keywords:} financial inclusion; access--usage gap; banking infrastructure; digital
payments; UPI; panel data; two-way fixed effects; India; pre-registration.
